\chapter{Inference and Model Compression}
\label{chap:inference}
Training asks how to learn a useful function.
Inference asks how to execute that learned function repeatedly under latency, throughput, memory and cost constraints.
The central design tension is
\[
\boxed{
\text{quality}
\longleftrightarrow
\text{latency}
\longleftrightarrow
\text{throughput}
\longleftrightarrow
\text{memory}
\longleftrightarrow
\text{cost}
}.
\]

\section{Quantisation: spend fewer bits on tolerant values}
\label{sec:inference-quantisation}
Quantisation maps continuous values onto a finite set of representable levels.
One useful scalar picture is
\[
\underbrace{q}_{\text{stored integer}}
=
\operatorname{round}\!\left(
\frac{\underbrace{x}_{\text{full-precision value}}}{\underbrace{s}_{\text{quantisation scale}}}
\right),
\qquad
\underbrace{\hat x}_{\text{reconstructed value}}
=sq.
\]
The approximation saves memory and often bandwidth/compute on suitable hardware, while introducing numeric error.
\figpdf[.76\textwidth]{quantisation_levels}{Quantisation replaces continuous values with nearby discrete levels; fewer levels increase approximation error but reduce representation cost.}

\subsection{PTQ versus QAT}
\label{subsec:inference-ptq-qat}
\textbf{Post-training quantisation (PTQ)} trains normally and quantises afterwards.
It is cheap and simple, but the trained model never adapted to the low-precision distortion.
\textbf{Quantisation-aware training (QAT)} simulates quantisation during training so the parameters can move toward a solution that survives deployment precision.

\begin{figure}[htbp]
\centering
\begin{tikzpicture}[node distance=7mm and 12mm]
\node[flow] (train1) {Train full\\precision};
\node[flow,right=of train1] (quant1) {Quantise\\afterwards};
\node[yesbox,right=of quant1] (ptq) {PTQ};
\node[flow,below=13mm of train1] (train2) {Train with simulated\\quantisation};
\node[flow,right=of train2] (quant2) {Deploy same\\low precision};
\node[yesbox,right=of quant2] (qat) {QAT};
\draw[arr] (train1)--(quant1);\draw[arr] (quant1)--(ptq);
\draw[arr] (train2)--(quant2);\draw[arr] (quant2)--(qat);
\end{tikzpicture}
\caption{PTQ tests whether a trained model survives quantisation; QAT exposes the model to that distortion while it can still adapt.}
\end{figure}

Activation-aware methods such as AWQ allocate extra care to weight directions whose quantisation error is more consequential for observed activations.
The broader principle is to spend precision where the function is sensitive instead of treating every parameter equally.

\begin{calloutmemoryhook}
PTQ tests whether the model survives low precision; QAT teaches it to survive.
\end{calloutmemoryhook}

\begin{calloutdegreeofexperiment}
Vary bit width, weight-only versus weight+activation quantisation, per-channel/group scaling, calibration data and which sensitive layers remain at higher precision.
\end{calloutdegreeofexperiment}

\section{Knowledge distillation: learn from a teacher distribution}
\label{sec:inference-distillation}
A hard label says only which class won.
A teacher distribution can also reveal which alternatives are plausible:
\[
[0.70,\ 0.20,\ 0.08,\ 0.02]
\quad\text{contains more relative structure than}\quad
[1,\ 0,\ 0,\ 0].
\]
A smaller student can therefore optimise both the ground-truth objective and a teacher-matching term:
\[
\underbrace{\mathcal L_{\text{student}}}_{\text{compact model objective}}
=
\underbrace{\lambda\mathcal L_{\text{labels}}}_{\text{task truth}}
+
\underbrace{(1-\lambda)\KL(P_{\text{teacher}}\|P_{\text{student}})}_{\text{match teacher structure}}.
\]
The student trades capacity for cheaper inference; the teacher's uncertainty/mistakes become part of the supervision.

\begin{calloutbias}
If the teacher's probability structure captures useful similarities between outputs, matching it can transfer more information than hard labels alone; teacher bias is transferred too.
\end{calloutbias}

\section{Cascades: reserve expensive models for hard cases}
\label{sec:inference-cascades}
If most inputs are easy, running the most expensive model on every request wastes compute.
A cascade routes easy/confident cases through a cheap model and escalates only uncertain or high-risk cases.

\begin{figure}[htbp]
\centering
\begin{tikzpicture}[node distance=8mm and 11mm]
\node[flow] (x) {Input};
\node[flow,right=of x] (small) {Cheap model};
\node[decisionnode,right=of small] (gate) {easy /\\confident?};
\node[yesbox,above right=5mm and 13mm of gate] (out1) {Return\\prediction};
\node[flow,below right=5mm and 13mm of gate] (large) {Expensive\\model};
\node[yesbox,right=of large] (out2) {Return\\prediction};
\draw[arr] (x)--(small);\draw[arr] (small)--(gate);
\draw[arr] (gate)--node[above,font=\scriptsize]{yes}(out1);
\draw[arr] (gate)--node[below,font=\scriptsize]{no}(large);
\draw[arr] (large)--(out2);
\end{tikzpicture}
\caption{A cascade turns model difficulty into a routing problem.}
\end{figure}

Expected compute becomes
\[
\underbrace{\E[C]}_{\text{average serving cost}}
=
\underbrace{C_{\text{cheap}}}_{\text{paid every request}}
+
\underbrace{P(\text{escalate})C_{\text{expensive}}}_{\text{paid only on routed cases}}.
\]
The gate is therefore as important as the component models.

\section{LLM inference: prefill versus decode}
\label{sec:inference-prefill-decode}
Autoregressive LLM serving contains two different phases.

\textbf{Prefill} processes the already-known prompt and creates the attention state for its tokens.
Prompt tokens can be processed with substantial parallelism because they are all available.
\textbf{Decode} generates one new token at a time because token $t+1$ cannot be chosen before token $t$ exists.

\begin{figure}[htbp]
\centering
\begin{tikzpicture}[node distance=7mm and 10mm]
\node[flow] (prompt) {Prompt tokens\\known together};
\node[flow,right=of prompt] (prefill) {Prefill\\parallel prompt work};
\node[yesbox,right=of prefill] (state) {KV state};
\node[flow,below=13mm of state] (decode) {Decode\\one new token};
\node[yesbox,right=of decode] (append) {Append token};
\draw[arr] (prompt)--(prefill);\draw[arr] (prefill)--(state);
\draw[arr] (state)--(decode);\draw[arr] (decode)--(append);
\draw[arr,rounded corners] (append.south)--++(0,-.6)--([yshift=-.6cm]decode.south)--(decode.south);
\end{tikzpicture}
\caption{Prefill builds state from a known prompt; decode repeatedly extends the sequence one token at a time.}
\end{figure}

\section{KV cache: save repeated attention work}
\label{sec:inference-kv-cache}
During autoregressive decode, old tokens keep the same key/value representations.
Recomputing them for every new token is wasteful, so serving systems cache
\[
(K_1,V_1),\ldots,(K_{N_{\text{context}}},V_{N_{\text{context}}}).
\]
The next token computes only its new query/key/value and reads the cached history.
This exchanges repeated compute for memory:
\[
\boxed{
\text{KV cache memory}
\propto
N_{\text{layers}}
\times
N_{\text{context tokens}}
\times
N_{\text{KV dims}}
\times
\text{bytes/value}
}.
\]
Long-context inference therefore becomes a memory-management problem even when model weights fit comfortably.

\begin{calloutmemoryhook}
KV cache remembers old attention keys/values so decode does not recompute the past.
\end{calloutmemoryhook}

\section{Batching and continuous batching}
\label{sec:inference-continuous-batching}
Static batching is efficient when requests have similar shapes.
Autoregressive requests finish at different output lengths, so a static batch can leave accelerator slots idle.
Continuous batching admits new requests as old requests finish.

\begin{center}
\small
\begin{tabularx}{.88\linewidth}{@{}L{.22\linewidth}Y@{}}
\toprule
\textbf{Static} & A: \code{============}\quad B: \code{===\phantom{=========}} \\
\textbf{Continuous} & A: \code{============}\quad B: \code{===}\quad C starts when B finishes \\
\bottomrule
\end{tabularx}
\end{center}

Paged KV-memory schemes similarly try to avoid reserving large contiguous cache blocks for every request.
The precise serving engine changes, but the systems goal is stable: keep expensive accelerator memory and compute occupied by useful work.

\section{Inference evolution and design axes}
\label{sec:inference-evolution}
\begin{center}
\small
\begin{tabularx}{.97\linewidth}{@{}L{.18\linewidth}L{.22\linewidth}YY@{}}
\toprule
\textbf{Step} & \textbf{What changed} & \textbf{Problem targeted} & \textbf{Trade-off} \\
\midrule
Full precision & Baseline execution & Maximum numeric fidelity & High memory/bandwidth \\
PTQ / QAT & Lower bit-width & Serving memory/throughput & Quantisation error / extra training \\
Distillation & Smaller student & Model size/latency & Student capacity limit \\
Cascade & Conditional model use & Waste on easy cases & Gate errors / system complexity \\
KV caching & Reuse old attention state & Repeated decode compute & Cache memory \\
Continuous batching & Dynamic request scheduling & Accelerator under-utilisation & More complex serving scheduler \\
\bottomrule
\end{tabularx}
\end{center}

\begin{calloutdegreeofexperiment}
Vary numeric precision, model size, cache representation, batch/scheduling policy, cascade gate, distillation objective and context budget.
\end{calloutdegreeofexperiment}




\section{Inference engines: the runtime around the model}
\label{sec:inference-engines}
A model defines which mathematical operations produce the next-token distribution.
An \textbf{inference engine} decides how those operations should actually execute under many concurrent requests, finite accelerator memory and latency constraints.

A useful separation is
\[
\boxed{
\text{model}
\rightarrow
\text{engine}
\rightarrow
\text{kernels / hardware}
}
\]
where the model defines the computation, the engine schedules and manages it, and kernels execute the resulting operations efficiently.

Running
\code{model.forward()}
answers only the first problem.
A serving engine must additionally decide which requests run now, how many tokens each receives, where their KV state lives, which kernels execute them and how work is divided across devices.

\begin{figure}[htbp]
\centering
\begin{tikzpicture}[node distance=7mm and 9mm]
\node[flow] (request) {Requests};
\node[flow,right=of request] (queue) {Admission /\\request queue};
\node[flow,right=of queue] (sched) {Scheduler};
\node[flow,right=of sched] (batch) {Physical\\batch};
\node[flow,below=11mm of batch] (runner) {Model runner};
\node[flow,left=of runner] (kernels) {GPU kernels};
\node[yesbox,left=of kernels] (token) {Sampled\\tokens};
\node[flow,below=11mm of sched] (kv) {KV-cache\\manager};

\draw[arr] (request)--(queue);
\draw[arr] (queue)--(sched);
\draw[arr] (sched)--(batch);
\draw[arr] (batch)--(runner);
\draw[arr] (runner)--(kernels);
\draw[arr] (kernels)--(token);
\draw[arr] (sched)--(kv);
\draw[arr] (kv)--(runner);
\draw[arr,rounded corners] (token.south)--++(0,-.7)--([yshift=-.7cm]sched.south)--(sched.south);
\end{tikzpicture}
\caption{An autoregressive inference engine repeatedly converts waiting requests and stored state into the next useful batch of GPU work.}
\label{fig:inference-engine-loop}
\end{figure}

The scheduler therefore runs inside a loop rather than constructing one fixed batch and executing it to completion.
After each iteration, some requests finish, some need another token, new requests arrive and KV memory becomes free or occupied.

A simplified engine iteration is
\[
\underbrace{
\text{waiting requests}
+
\text{active requests}
+
\text{available KV memory}
}_{\text{current serving state}}
\longrightarrow
\underbrace{\text{schedule next token work}}_{\text{engine decision}}
\longrightarrow
\underbrace{\text{GPU execution}}_{\text{model computation}}
\longrightarrow
\text{updated serving state}.
\]

Representative modern engines include \textbf{vLLM}, \textbf{SGLang} and \textbf{TensorRT-LLM}.
Their implementations differ, but they solve largely the same resource-allocation problem.

\begin{calloutmemoryhook}
The model defines the maths; the inference engine decides how that maths occupies the machine.
\end{calloutmemoryhook}


\section{Paged KV memory: variable sequences need an allocator}
\label{sec:inference-paged-kv}
Section~\ref{sec:inference-kv-cache} treated KV cache as saved attention state.
An engine must additionally decide where that state physically lives.

Suppose each request reserves one large contiguous KV region for its maximum possible sequence length.
A short generation then leaves much of that reservation unused, while variable sequence lengths make allocation and compaction awkward.

\begin{center}
\small
\begin{tabularx}{.96\linewidth}{@{}L{.22\linewidth}YY@{}}
\toprule
\textbf{Layout} & \textbf{Mechanism} & \textbf{Consequence} \\
\midrule
Large contiguous reservation
& Reserve one continuous region for each sequence
& Simple addressing, but variable lengths create wasted or fragmented memory \\
Paged / blocked KV
& Divide KV state into fixed-size blocks and map logical sequence positions onto physical blocks
& Easier allocation and reuse, at the cost of indirection and metadata \\
\bottomrule
\end{tabularx}
\end{center}

Paged KV management applies the same basic idea as virtual memory:
\[
\underbrace{\text{logical token positions}}_{\text{request sees one sequence}}
\longrightarrow
\underbrace{\text{block table}}_{\text{address mapping}}
\longrightarrow
\underbrace{\text{physical KV blocks}}_{\text{possibly non-contiguous GPU memory}}.
\]

\begin{figure}[htbp]
\centering
\begin{tikzpicture}[node distance=5mm and 9mm]
\node[flow] (seq) {Logical sequence\\tokens 1--10};
\node[flow,right=of seq] (map) {Block table};

\node[yesbox,right=12mm of map,yshift=9mm] (b1) {KV block 7};
\node[yesbox,right=12mm of map] (b2) {KV block 2};
\node[yesbox,right=12mm of map,yshift=-9mm] (b3) {KV block 11};

\draw[arr] (seq)--(map);
\draw[arr] (map)--(b1);
\draw[arr] (map)--(b2);
\draw[arr] (map)--(b3);
\end{tikzpicture}
\caption{Paged KV memory separates the logical sequence from its physical GPU-memory locations, so a growing sequence need not own one large contiguous reservation.}
\label{fig:paged-kv}
\end{figure}

This makes cache capacity a scheduling resource.
An engine can admit another request only if enough KV blocks are available for the work it expects to execute.

When memory pressure becomes too high, possible responses include delaying admission, evicting reusable state, preempting an active sequence, recomputing discarded state or moving state to a slower memory tier.
These choices exchange memory capacity against latency and repeated computation.

\begin{calloutdegreeofexperiment}
Vary KV block size, cache dtype, memory fraction assigned to KV state, eviction/preemption policy, maximum active sequences and whether cache state may spill to CPU or remote storage.
\end{calloutdegreeofexperiment}


\section{Prefix caching: identical history means identical old KV state}
\label{sec:inference-prefix-cache}
KV caching normally reuses history \emph{within} one request.
\textbf{Prefix caching} extends the same argument across requests.

If two requests begin with identical token prefixes under the same model state, their KV representations for that shared prefix are also identical.
The second request can therefore reference existing KV blocks instead of recomputing that part of prefill.

\[
\underbrace{
[\text{shared system prompt}\mid\text{shared document}]
}_{\text{existing cached prefix}}
+
\underbrace{[\text{new user suffix}]}_{\text{new work}}
\quad\Rightarrow\quad
\text{prefill only what is not reusable}.
\]

This is especially useful when many requests share large system prompts, retrieved documents, few-shot examples or previous conversation state.
It provides little benefit when prefixes rarely repeat.

The cache therefore introduces a second question alongside ordinary load:
\[
\boxed{
\text{best worker}
\neq
\text{least busy worker}
}
\]
because a somewhat busier worker holding most of the required prefix may require less total computation than an idle worker with no matching cache.

\begin{calloutbias}
If prompts contain repeated prefixes, cached KV state is reusable work; if prompts are mostly unique, maintaining and routing around that cache buys much less.
\end{calloutbias}


\section{Scheduling is the central inference-engine problem}
\label{sec:inference-scheduling}
Continuous batching from Section~\ref{sec:inference-continuous-batching} says that requests may enter and leave the active batch dynamically.
The engine still needs a policy for deciding \emph{which} work enters each iteration.

A useful abstraction is
\[
\underbrace{B_{\text{token budget}}}_{\text{GPU work allowed this iteration}}
=
\underbrace{N_{\text{decode tokens}}}_{\text{usually one per active decoder}}
+
\underbrace{N_{\text{prefill tokens}}}_{\text{new prompt work}}.
\]

The scheduler allocates this finite token budget across competing requests.
Larger batches often improve throughput, but waiting to construct them can increase latency.
Serving therefore turns batching into a repeated online optimisation problem rather than a fixed training hyperparameter.

\subsection{Chunked prefill}
\label{subsec:inference-chunked-prefill}
A long prompt can occupy substantial GPU time if its whole prefill executes in one iteration.
Existing decoding requests then wait longer for their next token.

\textbf{Chunked prefill} divides a long prompt into smaller pieces that can be interleaved with decode work:
\[
\text{long prefill}
=
\underbrace{\text{chunk}_1}_{\text{run}}
+
\underbrace{\text{chunk}_2}_{\text{later}}
+
\cdots.
\]

\begin{center}
\small
\begin{tabularx}{.96\linewidth}{@{}L{.23\linewidth}YY@{}}
\toprule
\textbf{Choice} & \textbf{Gain} & \textbf{Cost} \\
\midrule
Large prefill chunks
& Efficient large GPU work; prompt completes in fewer scheduling rounds
& Longer interruptions to active decode requests \\
Small prefill chunks
& Better opportunity to protect token-to-token latency
& More scheduling/launch fragmentation and possibly lower prefill efficiency \\
Separate prefill workers
& Stronger isolation between prefill and decode
& KV transfer and additional distributed-system complexity \\
\bottomrule
\end{tabularx}
\end{center}

The broader scheduler may additionally vary request priority, fairness, maximum active sequences, preemption policy and admission control.
These choices determine whether a saturated service fails through large queues, high tail latency or rejected work.

\begin{calloutmemoryhook}
Inference scheduling allocates scarce token work and KV memory across requests that arrive and finish at different times.
\end{calloutmemoryhook}

\begin{calloutdegreeofexperiment}
Vary maximum batched tokens, maximum active sequences, prefill chunk size, scheduling policy, request priority, admission limits and the balance between prefill and decode work.
\end{calloutdegreeofexperiment}


\section{FlashAttention and PagedAttention solve different problems}
\label{sec:inference-flash-paged}
Two similarly named ideas operate at different levels.

\textbf{FlashAttention} changes how an attention calculation moves intermediate data through the GPU memory hierarchy.
It avoids materialising unnecessarily large intermediate attention matrices in slow accelerator memory and therefore targets the efficiency of an attention operation.

\textbf{PagedAttention-style serving} changes how persistent KV state from many sequences is allocated and addressed.
It therefore targets serving-time memory management.

\[
\boxed{
\begin{array}{rcl}
\text{FlashAttention} &\rightarrow& \text{execute attention efficiently},\\[2mm]
\text{paged KV attention} &\rightarrow& \text{store and address many sequence states efficiently}.
\end{array}
}
\]

An engine can use both because they optimise different bottlenecks.
This distinction is useful more generally: \textbf{scheduler}, \textbf{memory manager} and \textbf{kernel} optimisations should not be treated as interchangeable.

\begin{calloutmemoryhook}
FlashAttention optimises attention computation; paged attention optimises serving-state memory.
\end{calloutmemoryhook}


\section{Below the scheduler: kernels, fusion and graphs}
\label{sec:inference-kernels}
Once the scheduler has constructed a batch, that logical work must become efficient GPU operations.
The execution path is approximately
\[
\text{scheduled tokens}
\rightarrow
\text{tensor shapes/layouts}
\rightarrow
\text{kernel selection}
\rightarrow
\text{GPU kernels}
\rightarrow
\text{outputs}.
\]

A Transformer block that conceptually contains normalisation, projections, RoPE, attention and an MLP may therefore execute through specialised kernels rather than a literal sequence of high-level framework operations.

\textbf{Kernel fusion} combines neighbouring operations when doing so avoids intermediate memory traffic or repeated launch overhead.
The benefit depends on the operations, tensor shapes and hardware; fusion that creates excessive register or shared-memory pressure can instead hurt performance.

\textbf{Compiled execution} specialises parts of the model graph for expected operation patterns and shapes.
\textbf{CUDA graphs} can replay a captured sequence of GPU operations with less CPU launch overhead when the execution structure is sufficiently stable.

The engine therefore sits between model-level semantics and hardware-level execution:
\[
\underbrace{\text{PyTorch / model operations}}_{\text{what should happen}}
\longrightarrow
\underbrace{\text{engine + compiler decisions}}_{\text{how to execute}}
\longrightarrow
\underbrace{\text{CUDA / Triton / specialised kernels}}_{\text{machine work}}.
\]

Attention is not the only important kernel.
Large dense projections are dominated by matrix multiplication, while sparse Mixture-of-Experts models additionally require expert routing, grouped matrix operations and communication between devices.

\begin{calloutpracticaltip}
A GPU that is not fully utilised may be waiting on the CPU scheduler, kernel launches, communication or memory rather than lacking enough model FLOPs.
\end{calloutpracticaltip}


\section{Speculative decoding: predict cheaply, verify expensively}
\label{sec:inference-speculative}
Standard autoregressive decode requires one target-model step before each accepted token.
This serial dependency limits how much token generation can be parallelised.

\textbf{Speculative decoding} uses a cheaper mechanism to propose several future tokens and then asks the target model to verify those proposals together.

\begin{figure}[htbp]
\centering
\begin{tikzpicture}[node distance=7mm and 10mm]
\node[flow] (state) {Current\\sequence};
\node[flow,right=of state] (draft) {Cheap proposer\\predicts several};
\node[flow,right=of draft] (verify) {Target model\\verifies batch};
\node[decisionnode,right=of verify] (accepted) {accepted\\prefix};
\node[yesbox,right=of accepted] (out) {Append accepted\\tokens};

\draw[arr] (state)--(draft);
\draw[arr] (draft)--(verify);
\draw[arr] (verify)--(accepted);
\draw[arr] (accepted)--(out);
\draw[arr,rounded corners] (out.south)--++(0,-.7)--([yshift=-.7cm]state.south)--(state.south);
\end{tikzpicture}
\caption{Speculative decoding spends cheap work proposing multiple tokens so that one target-model verification can sometimes advance the sequence by more than one token.}
\label{fig:speculative-decoding}
\end{figure}

If the proposer predicts the target distribution well, several tokens are accepted per expensive target-model pass.
If proposals are poor, the system pays proposal and verification overhead while advancing by few additional tokens.

A useful quantity is therefore
\[
\underbrace{N_{\text{accepted per verification}}}_{\text{speculative progress}}
\]
rather than merely the number of tokens proposed.

Possible proposers include a smaller draft model, simple continuation rules, extra prediction heads or specialised learned draft mechanisms.
The common assumption is the same: future tokens can be predicted cheaply enough and accurately enough that verification saves target-model serial work.

\begin{calloutbias}
If cheap proposals strongly match the target model, speculative decoding converts prediction accuracy into lower serial decode cost; poor acceptance can make the extra work useless.
\end{calloutbias}

\begin{calloutdegreeofexperiment}
Vary proposer architecture, number of drafted tokens, acceptance rule, draft/target size ratio and when speculation is enabled for a request.
\end{calloutdegreeofexperiment}


\section{Distributed inference: divide weights, layers, experts, context or requests}
\label{sec:inference-distributed}
One accelerator may be insufficient because the model does not fit, the KV cache consumes too much memory or one replica cannot provide the required throughput.

Different forms of parallelism divide different things.

\begin{center}
\small
\begin{tabularx}{.98\linewidth}{@{}L{.16\linewidth}L{.20\linewidth}YY@{}}
\toprule
\textbf{Parallelism} & \textbf{What is divided?} & \textbf{Useful pressure} & \textbf{Main new cost} \\
\midrule
Data
& Requests / batches
& More aggregate serving throughput
& Replicated model weights; routing/load balance \\
Tensor
& Tensor operations within layers
& Individual layers/model state need several GPUs
& Frequent communication inside layers \\
Pipeline
& Groups of layers
& Model spans devices
& Pipeline bubbles and stage imbalance \\
Expert
& MoE experts
& Sparse expert parameters distributed across GPUs
& Token routing / all-to-all communication \\
Context
& Sequence positions
& Very long contexts or attention state
& Communication around distributed attention/state \\
\bottomrule
\end{tabularx}
\end{center}

Parallelism is therefore not simply ``use more GPUs''.
It changes both memory placement and the communication graph.

For example, tensor parallelism can make a model fit by dividing weights across devices, but every token may then require communication between those devices.
Adding GPUs can therefore lower latency when it removes a compute/memory bottleneck or increase latency when communication dominates.

\[
\boxed{
\text{more devices}
\rightarrow
\text{more aggregate compute/memory}
\quad\longleftrightarrow\quad
\text{more communication}
}
\]

The best serving topology depends on model architecture, request lengths, interconnect bandwidth, latency target and whether the objective is per-request speed or aggregate throughput.

\begin{calloutdegreeofexperiment}
Vary tensor-, pipeline-, data-, expert- and context-parallel sizes, replica count, device placement and the mapping between parallel ranks and physical interconnect topology.
\end{calloutdegreeofexperiment}


\section{Disaggregated serving: prefill and decode need not share GPUs}
\label{sec:inference-disaggregation}
Section~\ref{sec:inference-prefill-decode} separated autoregressive inference into prefill and decode mathematically.
A modern serving system can also separate them \emph{physically}.

Prefill typically provides large amounts of parallel token work, while decode repeatedly operates on small new-token batches and a growing KV cache.
One shared worker must compromise between these different execution profiles.

\textbf{Prefill/decode disaggregation} instead uses different accelerator pools:
\[
\underbrace{\text{request}}_{\text{prompt}}
\rightarrow
\underbrace{\text{prefill worker}}_{\text{build KV state}}
\rightarrow
\underbrace{\text{KV transfer}}_{\text{move state}}
\rightarrow
\underbrace{\text{decode worker}}_{\text{generate tokens}}.
\]

\begin{figure}[htbp]
\centering
\begin{tikzpicture}[node distance=8mm and 12mm]
\node[flow] (req) {Request};
\node[flow,right=of req] (router) {Router};
\node[flow,right=of router] (prefill) {Prefill\\GPU pool};
\node[yesbox,right=of prefill] (kv) {KV cache};
\node[flow,right=of kv] (decode) {Decode\\GPU pool};
\node[yesbox,right=of decode] (out) {Generated\\tokens};

\draw[arr] (req)--(router);
\draw[arr] (router)--(prefill);
\draw[arr] (prefill)--(kv);
\draw[arr] (kv)--node[above,font=\scriptsize]{GPU / network transfer}(decode);
\draw[arr] (decode)--(out);
\end{tikzpicture}
\caption{Disaggregated serving allows prefill and decode to use independently sized and configured worker pools, but KV state must cross the boundary.}
\label{fig:prefill-decode-disaggregation}
\end{figure}

The benefit is independent control.
Prefill and decode pools may use different replica counts, parallelism strategies or even hardware chosen for their respective bottlenecks.

The cost is that potentially large KV state must move between workers before decoding can continue.
Fast GPU-to-GPU links and overlap between communication and computation therefore become part of the inference-engine problem.

Disaggregation is not automatically faster.
For short prompts, low concurrency or slow interconnects, keeping prefill and decode together can avoid transfer and routing overhead.

\begin{calloutbias}
If prefill and decode have sufficiently different resource requirements, separating them allows independent optimisation; the price is moving KV state between them.
\end{calloutbias}

\begin{calloutdegreeofexperiment}
Vary prefill/decode worker ratio, parallelism within each pool, KV-transfer mechanism, placement topology, prefill chunking and whether requests are selectively or always disaggregated.
\end{calloutdegreeofexperiment}


\section{Routing becomes cache-aware}
\label{sec:inference-kv-routing}
Once a service contains many workers, a frontend must choose where each request executes.

Simple round-robin routing assumes workers are interchangeable.
KV caching breaks that assumption because different workers may already hold different reusable prefixes.

A cache-aware routing decision therefore balances at least
\[
\boxed{
\text{routing cost}
\approx
\underbrace{\text{new computation required}}_{\text{prefix not cached}}
+
\underbrace{\text{existing worker load}}_{\text{queue / active requests}}
+
\underbrace{\text{transfer cost}}_{\text{if state must move}}
}.
\]

Sending every repeated prompt to its cache-rich worker maximises reuse but can create a hot worker.
Sending everything to the emptiest worker balances load but can repeatedly throw away useful cached computation.
The router therefore trades \textbf{locality} against \textbf{load}.

This is the point where an inference \emph{engine} becomes part of a larger inference \emph{system}.
Systems such as NVIDIA Dynamo can sit above engine workers and coordinate request routing, worker discovery, KV-aware placement and prefill/decode pools.

\begin{calloutmemoryhook}
With reusable KV state, routing decides where the computation is cheapest, not merely where the queue is shortest.
\end{calloutmemoryhook}


\section{Inference metrics: latency has internal structure}
\label{sec:inference-metrics}
A single average request latency hides different serving failures.
Autoregressive systems commonly separate the user-visible path into
\[
\underbrace{\text{TTFT}}_{\text{request}\rightarrow\text{first token}}
\qquad\text{and}\qquad
\underbrace{\text{TPOT / ITL}}_{\text{spacing between later tokens}}.
\]

\textbf{Time to first token (TTFT)} contains queueing plus prompt/prefill work before useful output begins.
\textbf{Inter-token latency (ITL)} measures the gaps between streamed tokens, while \textbf{time per output token (TPOT)} summarises generation time after the first token.
Throughput reports completed requests or tokens per unit time.

These metrics can move in opposite directions.
Waiting briefly to build a larger batch may increase throughput while worsening latency; reserving resources for low-latency decode may improve ITL while reducing maximum aggregate throughput.

The workload must therefore be specified alongside a benchmark:
\[
\boxed{
\text{performance}
=
f(
N_{\text{input tokens}},
N_{\text{output tokens}},
\text{arrival rate},
\text{concurrency},
\text{hardware},
\text{precision},
\text{parallelism},
\text{cache reuse}
).
}
\]

``Tokens/s'' without those conditions is not a complete comparison.

Tail metrics matter because overloaded queues, long prompts and unlucky scheduling decisions may affect a minority of requests far more than the mean.
A service with acceptable mean TTFT but unacceptable p99 TTFT still provides poor latency guarantees to some users.

\begin{center}
\small
\begin{tabularx}{.98\linewidth}{@{}L{.18\linewidth}L{.30\linewidth}Y@{}}
\toprule
\textbf{Symptom} & \textbf{Likely pressure} & \textbf{Questions to ask} \\
\midrule
High TTFT
& Queueing or expensive prefill
& Is arrival rate above capacity; are prompts long; would prefix reuse/chunking help? \\
High ITL / TPOT
& Decode scheduling, memory bandwidth or communication
& Are prefills interrupting decode; is the active batch too large; is distributed communication expensive? \\
Low throughput
& GPU under-utilisation or inefficient execution
& Is the GPU waiting on CPU scheduling, kernels, communication or too-small batches? \\
KV OOM / preemption
& Too much active sequence state
& Is context/concurrency too high; can KV precision, allocation or admission policy change? \\
Good mean, poor p99
& Queueing or heterogeneous request interference
& Which request class causes stalls; should admission, priority or P/D isolation change? \\
\bottomrule
\end{tabularx}
\end{center}

\begin{calloutpracticaltip}
Benchmark the request-length and concurrency distribution you expect to serve; engine rankings can reverse when the workload changes.
\end{calloutpracticaltip}


\section{Modern inference-engine design space}
\label{sec:inference-engine-design-space}
Modern engines combine optimisations at several distinct layers.

\begin{center}
\small
\begin{tabularx}{.98\linewidth}{@{}L{.18\linewidth}L{.23\linewidth}YY@{}}
\toprule
\textbf{Layer} & \textbf{Example choice} & \textbf{What becomes cheaper} & \textbf{New cost / constraint} \\
\midrule
Numeric representation
& Quantisation
& Weight/KV memory and supported low-precision compute
& Approximation error \\
Request scheduling
& Continuous batching / chunked prefill
& Idle accelerator time and request interference
& Scheduler complexity / latency trade-offs \\
State memory
& Paged KV allocation
& Variable-length cache allocation
& Address indirection / metadata \\
Cross-request reuse
& Prefix caching
& Repeated prefill
& Cache capacity, lookup and placement complexity \\
Execution
& Fused kernels / CUDA graphs / compilation
& Memory traffic and launch overhead
& Specialisation and implementation complexity \\
Decode algorithm
& Speculative decoding
& Serial target-model steps
& Proposal + verification work \\
Device scaling
& Tensor / pipeline / expert parallelism
& Per-device model/compute requirement
& Communication \\
Fleet architecture
& P/D disaggregation
& Resource interference between phases
& KV transfer and routing \\
Fleet routing
& KV-aware placement
& Recomputed shared prefixes
& Global state and load-balancing complexity \\
\bottomrule
\end{tabularx}
\end{center}

The important distinction is \emph{which bottleneck an optimisation attacks}.
Quantisation does not solve scheduling, paged KV allocation does not make an individual attention kernel fast, and a faster kernel does not prevent idle slots caused by poor request scheduling.

\begin{calloutdegreeofexperiment}
The full inference-engine experiment space includes numeric precision, KV representation, batch/token budget, scheduling policy, cache allocation, prefix reuse, kernel/backend choice, speculative method, parallelism topology, worker counts, routing policy and prefill/decode placement.
\end{calloutdegreeofexperiment}


\section{Inference-engine evolution}
\label{sec:inference-engine-evolution}
The serving stack evolved by repeatedly moving optimisation outward from one model execution toward the complete request fleet.

\begin{center}
\small
\begin{tabularx}{.98\linewidth}{@{}L{.18\linewidth}L{.23\linewidth}YY@{}}
\toprule
\textbf{Step} & \textbf{What changed} & \textbf{Pressure addressed} & \textbf{New trade-off} \\
\midrule
Naive generation
& Execute one request/token loop
& Correct model output
& Poor hardware utilisation \\
Static batching
& Execute several requests together
& Small GPU operations
& Wait for longest sequence / rigid batch \\
KV caching
& Preserve old attention state
& Repeated attention computation
& Growing memory use \\
Continuous batching
& Replace finished requests during generation
& Idle batch slots
& Online scheduling complexity \\
Paged KV
& Allocate sequence state in blocks
& Fragmentation / variable lengths
& Mapping metadata and allocator complexity \\
Prefix caching
& Reuse KV across requests
& Repeated shared prefill
& Cache management / routing locality \\
Chunked prefill
& Interleave portions of long prompts
& Decode interruption
& Prefill fragmentation \\
Speculative decoding
& Cheaply propose then verify several tokens
& Serial decode dependency
& Draft/verification overhead \\
P/D disaggregation
& Separate prompt and generation workers
& Different prefill/decode resource profiles
& KV communication \\
KV-aware fleet routing
& Route using cache state plus load
& Duplicate work across replicas
& Distributed metadata / coordination \\
\bottomrule
\end{tabularx}
\end{center}

The progression can be compressed to
\[
\boxed{
\begin{aligned}
\text{make one model execution efficient}
&\rightarrow
\text{keep the accelerator occupied}\\
&\rightarrow
\text{manage sequence state efficiently}\\
&\rightarrow
\text{avoid unnecessary computation}\\
&\rightarrow
\text{optimise placement across the fleet}.
\end{aligned}
}
\]

\begin{calloutmemoryhook}
Inference engines evolved from making a forward pass fast to deciding which useful work should occupy every byte, kernel and GPU.
\end{calloutmemoryhook}


\section{Engine families: same problem, different emphasis}
\label{sec:inference-engine-families}
Product details change quickly, so engine names are more useful as examples of design emphasis than as a fixed ranking.

\begin{center}
\small
\begin{tabularx}{.98\linewidth}{@{}L{.18\linewidth}L{.27\linewidth}YY@{}}
\toprule
\textbf{System} & \textbf{Useful mental model} & \textbf{Characteristic emphasis} & \textbf{Role} \\
\midrule
vLLM
& General high-throughput model runtime
& Paged KV management, flexible scheduling, broad model/kernel support
& Inference engine \\
SGLang
& Runtime where reusable request structure is explicit
& Prefix/KV reuse, scheduling and large-scale LLM serving
& Inference engine \\
TensorRT-LLM
& NVIDIA-specific model/runtime co-optimisation
& Specialised kernels, low precision, batching and distributed execution
& Inference engine \\
NVIDIA Dynamo
& Coordinate many engine workers as one serving system
& Routing, discovery, KV-aware placement and disaggregated serving
& Fleet / distributed-serving layer \\
\bottomrule
\end{tabularx}
\end{center}

Choosing between engines should therefore begin with the workload and hardware rather than a benchmark headline.
Relevant questions include model support, quantisation requirements, prefix-reuse structure, request lengths, latency target, accelerator type, distributed topology and which parts of the stack need to remain customisable.

The stable abstraction is
\[
\boxed{
\text{request}
\rightarrow
\text{route}
\rightarrow
\text{schedule}
\rightarrow
\text{allocate state}
\rightarrow
\text{execute kernels}
\rightarrow
\text{sample}
\rightarrow
\text{repeat}.
}
\]

Once this flow is understood, an unfamiliar engine can be inspected by asking where it makes each decision and which resource that decision is trying to save.